% --- Archon physics formalization source begin ---
% archon:physics
% archon:covers PhyXMiniProblems/problem_phyx_mini_0919.lean
% archon:source-report reports/phyx_mini/problem_phyx_mini_0919.source.json

\chapter{Physics problem phyx\_mini\_0919}
\label{ch:PhyXMiniProblems_problem_phyx_mini_0919}

\paragraph{Problem source.}
\#\# Physical scenario

In the slidewire, energy is dissipated in the circuit owing to its resistance. Let the resistance of the circuit (made up of the moving slidewire and the U-shaped conductor that connects the ends of the slidewire) at a given point in the slidewire's motion be \$R\$.

\#\# Question

Find the rate at which energy is dissipated in the circuit.

\#\# Answer choices

- A: A: \textbackslash{}frac\{BLv\}\{R\}

- B: B: \textbackslash{}frac\{B\textasciicircum{}2 L v\textasciicircum{}2\}\{R\}

- C: C: \textbackslash{}frac\{B\textasciicircum{}2 L\textasciicircum{}2 v\textasciicircum{}2\}\{R\}

- D: D: \textbackslash{}frac\{B\textasciicircum{}2 L\textasciicircum{}2 v\}\{R\textasciicircum{}2\}

\#\# Figure caption (auxiliary; use the image as primary evidence)

The image depicts a conductive rectangular loop with a sliding wire in the presence of a uniform magnetic field. 

- **Objects:**
  - The loop is shown in a light brown color with parallel sides to represent wires.
  - A vertical sliding wire is depicted in the middle of the loop.

- **Magnetic Field:**
  - A uniform magnetic field is indicated by blue crosses (×), meaning the field is directed into the page.

- **Forces and Directions:**
  - Several vectors and symbols represent different directions and forces:
    - \textbackslash{}(\textbackslash{}vec\{B\}\textbackslash{}): Magnetic field vector, directed into the plane.
    - \textbackslash{}(\textbackslash{}vec\{I\}\textbackslash{}): Current vectors, depicted with purple arrows along the wires, showing current flow both clockwise and counterclockwise.
    - \textbackslash{}(\textbackslash{}vec\{F\}\textbackslash{}): Force vector in red, directed left on the sliding wire.
    - \textbackslash{}(\textbackslash{}vec\{v\}\textbackslash{}): Velocity vector in green, pointing to the right on the sliding wire.
    - \textbackslash{}(\textbackslash{}vec\{L\}\textbackslash{}): Length vector in blue, indicating the vertical orientation of the sliding wire.
    - \textbackslash{}(\textbackslash{}mathcal\{E\}\textbackslash{}): An electromotive force is represented by a blue arrow pointing upwards.

- **Relationships:**
  - The sliding wire moves to the right, inducing an emf (\textbackslash{}(\textbackslash{}mathcal\{E\}\textbackslash{})) due to motion in the magnetic field.
  - The force (\textbackslash{}(\textbackslash{}vec\{F\}\textbackslash{})) opposes the motion of the sliding wire, consistent with Lenz's Law.

This diagram is often used to demonstrate electromagnetic induction, highlighting the interaction between induced current, magnetic fields, and motion.

\#\# Dataset metadata

Electromagnetic Induction; Physical Model Grounding Reasoning, Multi-Formula Reasoning

\paragraph{Recorded answer/context.}
C: C: \textbackslash{}frac\{B\textasciicircum{}2 L\textasciicircum{}2 v\textasciicircum{}2\}\{R\}

\paragraph{Figure/image path.}
/root/proposal\_for\_physic/science-mango/phyx\_mini\_run/phyx\_data/test\_image/919.png

\paragraph{Formalization target.}
create a compiling Lean file with sorry bodies at `PhyXMiniProblems/problem\_phyx\_mini\_0919.lean`. The Lean declarations must preserve the physical quantities, dimensions or dimensional roles, figure labels, governing-law hypotheses, and final relation expressed by this problem.
Use Mathlib/Physlib names found through LeanExplore where available. If a physics API is missing, introduce faithful local abstractions rather than scalar placeholder aliases.

\begin{theorem}[Physics formalization target]
\label{thm:physics:phyx_mini_0919:target}
\lean{PhyXMiniProblems.ProblemPhyXMini0919.slidewire_energy_dissipation_rate}
\uses{def:physics:phyx-mini-0919:phyxminiproblems-problemphyxmini0919-magneticfluxdensityinteslas, def:physics:phyx-mini-0919:phyxminiproblems-problemphyxmini0919-lengthinmeters, def:physics:phyx-mini-0919:phyxminiproblems-problemphyxmini0919-speedinmeterspersecond, def:physics:phyx-mini-0919:phyxminiproblems-problemphyxmini0919-resistanceinohms, def:physics:phyx-mini-0919:phyxminiproblems-problemphyxmini0919-powerinwatts, def:physics:phyx-mini-0919:phyxminiproblems-problemphyxmini0919-slidewirecircuitsetup, def:physics:phyx-mini-0919:phyxminiproblems-problemphyxmini0919-matchesslidewirecircuitscenario, def:physics:phyx-mini-0919:phyxminiproblems-problemphyxmini0919-matchesprimaryslidewirefigure, def:physics:phyx-mini-0919:phyxminiproblems-problemphyxmini0919-hasphysicalslidewireparameters, def:physics:phyx-mini-0919:phyxminiproblems-problemphyxmini0919-satisfiescircuitresistancecomposition, def:physics:phyx-mini-0919:phyxminiproblems-problemphyxmini0919-hasuniformappliedmagneticfield, def:physics:phyx-mini-0919:phyxminiproblems-problemphyxmini0919-satisfiesslidewireinductionandcircuitlaws}
The assigned autoformalize agent should translate this physics problem into Lean declarations in the covered file, with theorem and lemma proof bodies written as `by sorry`.
\end{theorem}
\begin{proof}
This is an autoformalization task, not a proof task. Produce statements that can later be proved by the physics prover without weakening the physical model.
\end{proof}

% --- Archon declaration topology begin ---
\section{Declaration topology}
\begin{definition}[magnetic Flux Density Dimension]
  \label{def:physics:phyx-mini-0919:phyxminiproblems-problemphyxmini0919-magneticfluxdensitydimension}
  \lean{PhyXMiniProblems.ProblemPhyXMini0919.magneticFluxDensityDimension}
  The dimension M T⁻¹ C⁻¹ of magnetic flux density (tesla)
\end{definition}
\begin{proof}
  This is the declared model definition of magnetic Flux Density Dimension.
\end{proof}
\begin{definition}[electromotive Force Dimension]
  \label{def:physics:phyx-mini-0919:phyxminiproblems-problemphyxmini0919-electromotiveforcedimension}
  \lean{PhyXMiniProblems.ProblemPhyXMini0919.electromotiveForceDimension}
  The dimension M L² T⁻² C⁻¹ of electromotive force (volt)
\end{definition}
\begin{proof}
  This is the declared model definition of electromotive Force Dimension.
\end{proof}
\begin{definition}[electric Current Dimension]
  \label{def:physics:phyx-mini-0919:phyxminiproblems-problemphyxmini0919-electriccurrentdimension}
  \lean{PhyXMiniProblems.ProblemPhyXMini0919.electricCurrentDimension}
  The dimension C T⁻¹ of electric current (ampere)
\end{definition}
\begin{proof}
  This is the declared model definition of electric Current Dimension.
\end{proof}
\begin{definition}[electrical Resistance Dimension]
  \label{def:physics:phyx-mini-0919:phyxminiproblems-problemphyxmini0919-electricalresistancedimension}
  \lean{PhyXMiniProblems.ProblemPhyXMini0919.electricalResistanceDimension}
  The dimension M L² T⁻¹ C⁻² of electrical resistance (ohm)
\end{definition}
\begin{proof}
  This is the declared model definition of electrical Resistance Dimension.
\end{proof}
\begin{definition}[power Dimension]
  \label{def:physics:phyx-mini-0919:phyxminiproblems-problemphyxmini0919-powerdimension}
  \lean{PhyXMiniProblems.ProblemPhyXMini0919.powerDimension}
  The dimension M L² T⁻³ of power (watt)
\end{definition}
\begin{proof}
  This is the declared model definition of power Dimension.
\end{proof}
\begin{definition}[force Dimension]
  \label{def:physics:phyx-mini-0919:phyxminiproblems-problemphyxmini0919-forcedimension}
  \lean{PhyXMiniProblems.ProblemPhyXMini0919.forceDimension}
  The dimension M L T⁻² of force (newton)
\end{definition}
\begin{proof}
  This is the declared model definition of force Dimension.
\end{proof}
\begin{definition}[Magnetic Flux Density Magnitude]
  \label{def:physics:phyx-mini-0919:phyxminiproblems-problemphyxmini0919-magneticfluxdensitymagnitude}
  \lean{PhyXMiniProblems.ProblemPhyXMini0919.MagneticFluxDensityMagnitude}
  \uses{def:physics:phyx-mini-0919:phyxminiproblems-problemphyxmini0919-magneticfluxdensitydimension}
  A nonnegative, unit-independent magnetic-flux-density magnitude
\end{definition}
\begin{proof}
  This is the declared model definition of Magnetic Flux Density Magnitude.
\end{proof}
\begin{definition}[Length Magnitude]
  \label{def:physics:phyx-mini-0919:phyxminiproblems-problemphyxmini0919-lengthmagnitude}
  \lean{PhyXMiniProblems.ProblemPhyXMini0919.LengthMagnitude}
  A nonnegative, unit-independent physical length
\end{definition}
\begin{proof}
  This is the declared model definition of Length Magnitude.
\end{proof}
\begin{definition}[Speed Magnitude]
  \label{def:physics:phyx-mini-0919:phyxminiproblems-problemphyxmini0919-speedmagnitude}
  \lean{PhyXMiniProblems.ProblemPhyXMini0919.SpeedMagnitude}
  A nonnegative, unit-independent speed
\end{definition}
\begin{proof}
  This is the declared model definition of Speed Magnitude.
\end{proof}
\begin{definition}[Electromotive Force Magnitude]
  \label{def:physics:phyx-mini-0919:phyxminiproblems-problemphyxmini0919-electromotiveforcemagnitude}
  \lean{PhyXMiniProblems.ProblemPhyXMini0919.ElectromotiveForceMagnitude}
  \uses{def:physics:phyx-mini-0919:phyxminiproblems-problemphyxmini0919-electromotiveforcedimension}
  A nonnegative, unit-independent electromotive-force magnitude
\end{definition}
\begin{proof}
  This is the declared model definition of Electromotive Force Magnitude.
\end{proof}
\begin{definition}[Electric Current Magnitude]
  \label{def:physics:phyx-mini-0919:phyxminiproblems-problemphyxmini0919-electriccurrentmagnitude}
  \lean{PhyXMiniProblems.ProblemPhyXMini0919.ElectricCurrentMagnitude}
  \uses{def:physics:phyx-mini-0919:phyxminiproblems-problemphyxmini0919-electriccurrentdimension}
  A nonnegative, unit-independent electric-current magnitude
\end{definition}
\begin{proof}
  This is the declared model definition of Electric Current Magnitude.
\end{proof}
\begin{definition}[Resistance Magnitude]
  \label{def:physics:phyx-mini-0919:phyxminiproblems-problemphyxmini0919-resistancemagnitude}
  \lean{PhyXMiniProblems.ProblemPhyXMini0919.ResistanceMagnitude}
  \uses{def:physics:phyx-mini-0919:phyxminiproblems-problemphyxmini0919-electricalresistancedimension}
  A nonnegative, unit-independent electrical resistance
\end{definition}
\begin{proof}
  This is the declared model definition of Resistance Magnitude.
\end{proof}
\begin{definition}[Power Magnitude]
  \label{def:physics:phyx-mini-0919:phyxminiproblems-problemphyxmini0919-powermagnitude}
  \lean{PhyXMiniProblems.ProblemPhyXMini0919.PowerMagnitude}
  \uses{def:physics:phyx-mini-0919:phyxminiproblems-problemphyxmini0919-powerdimension}
  A nonnegative, unit-independent rate of energy dissipation
\end{definition}
\begin{proof}
  This is the declared model definition of Power Magnitude.
\end{proof}
\begin{definition}[Force Magnitude]
  \label{def:physics:phyx-mini-0919:phyxminiproblems-problemphyxmini0919-forcemagnitude}
  \lean{PhyXMiniProblems.ProblemPhyXMini0919.ForceMagnitude}
  \uses{def:physics:phyx-mini-0919:phyxminiproblems-problemphyxmini0919-forcedimension}
  A nonnegative, unit-independent magnetic-force magnitude
\end{definition}
\begin{proof}
  This is the declared model definition of Force Magnitude.
\end{proof}
\begin{definition}[magnetic Flux Density In Teslas]
  \label{def:physics:phyx-mini-0919:phyxminiproblems-problemphyxmini0919-magneticfluxdensityinteslas}
  \lean{PhyXMiniProblems.ProblemPhyXMini0919.magneticFluxDensityInTeslas}
  \uses{def:physics:phyx-mini-0919:phyxminiproblems-problemphyxmini0919-magneticfluxdensitymagnitude}
  Coherent-SI readout of magnetic flux density, in teslas
\end{definition}
\begin{proof}
  This is the declared model definition of magnetic Flux Density In Teslas.
\end{proof}
\begin{definition}[length In Meters]
  \label{def:physics:phyx-mini-0919:phyxminiproblems-problemphyxmini0919-lengthinmeters}
  \lean{PhyXMiniProblems.ProblemPhyXMini0919.lengthInMeters}
  \uses{def:physics:phyx-mini-0919:phyxminiproblems-problemphyxmini0919-lengthmagnitude}
  Coherent-SI readout of length, in metres
\end{definition}
\begin{proof}
  This is the declared model definition of length In Meters.
\end{proof}
\begin{definition}[speed In Meters Per Second]
  \label{def:physics:phyx-mini-0919:phyxminiproblems-problemphyxmini0919-speedinmeterspersecond}
  \lean{PhyXMiniProblems.ProblemPhyXMini0919.speedInMetersPerSecond}
  \uses{def:physics:phyx-mini-0919:phyxminiproblems-problemphyxmini0919-speedmagnitude}
  Coherent-SI readout of speed, in metres per second
\end{definition}
\begin{proof}
  This is the declared model definition of speed In Meters Per Second.
\end{proof}
\begin{definition}[electromotive Force In Volts]
  \label{def:physics:phyx-mini-0919:phyxminiproblems-problemphyxmini0919-electromotiveforceinvolts}
  \lean{PhyXMiniProblems.ProblemPhyXMini0919.electromotiveForceInVolts}
  \uses{def:physics:phyx-mini-0919:phyxminiproblems-problemphyxmini0919-electromotiveforcemagnitude}
  Coherent-SI readout of electromotive force, in volts
\end{definition}
\begin{proof}
  This is the declared model definition of electromotive Force In Volts.
\end{proof}
\begin{definition}[electric Current In Amperes]
  \label{def:physics:phyx-mini-0919:phyxminiproblems-problemphyxmini0919-electriccurrentinamperes}
  \lean{PhyXMiniProblems.ProblemPhyXMini0919.electricCurrentInAmperes}
  \uses{def:physics:phyx-mini-0919:phyxminiproblems-problemphyxmini0919-electriccurrentmagnitude}
  Coherent-SI readout of electric current, in amperes
\end{definition}
\begin{proof}
  This is the declared model definition of electric Current In Amperes.
\end{proof}
\begin{definition}[resistance In Ohms]
  \label{def:physics:phyx-mini-0919:phyxminiproblems-problemphyxmini0919-resistanceinohms}
  \lean{PhyXMiniProblems.ProblemPhyXMini0919.resistanceInOhms}
  \uses{def:physics:phyx-mini-0919:phyxminiproblems-problemphyxmini0919-resistancemagnitude}
  Coherent-SI readout of resistance, in ohms
\end{definition}
\begin{proof}
  This is the declared model definition of resistance In Ohms.
\end{proof}
\begin{definition}[power In Watts]
  \label{def:physics:phyx-mini-0919:phyxminiproblems-problemphyxmini0919-powerinwatts}
  \lean{PhyXMiniProblems.ProblemPhyXMini0919.powerInWatts}
  \uses{def:physics:phyx-mini-0919:phyxminiproblems-problemphyxmini0919-powermagnitude}
  Coherent-SI readout of dissipated power, in watts
\end{definition}
\begin{proof}
  This is the declared model definition of power In Watts.
\end{proof}
\begin{definition}[force In Newtons]
  \label{def:physics:phyx-mini-0919:phyxminiproblems-problemphyxmini0919-forceinnewtons}
  \lean{PhyXMiniProblems.ProblemPhyXMini0919.forceInNewtons}
  \uses{def:physics:phyx-mini-0919:phyxminiproblems-problemphyxmini0919-forcemagnitude}
  Coherent-SI readout of force, in newtons
\end{definition}
\begin{proof}
  This is the declared model definition of force In Newtons.
\end{proof}
\begin{definition}[Conductor Part]
  \label{def:physics:phyx-mini-0919:phyxminiproblems-problemphyxmini0919-conductorpart}
  \lean{PhyXMiniProblems.ProblemPhyXMini0919.ConductorPart}
  The two conducting parts whose resistances make up the circuit resistance
\end{definition}
\begin{proof}
  This is the declared model definition of Conductor Part.
\end{proof}
\begin{definition}[Spatial Direction]
  \label{def:physics:phyx-mini-0919:phyxminiproblems-problemphyxmini0919-spatialdirection}
  \lean{PhyXMiniProblems.ProblemPhyXMini0919.SpatialDirection}
  Named directions needed to transcribe the supplied side-view figure
\end{definition}
\begin{proof}
  This is the declared model definition of Spatial Direction.
\end{proof}
\begin{definition}[Circuit Sense]
  \label{def:physics:phyx-mini-0919:phyxminiproblems-problemphyxmini0919-circuitsense}
  \lean{PhyXMiniProblems.ProblemPhyXMini0919.CircuitSense}
  Orientation of current around the rectangular circuit as seen on the page
\end{definition}
\begin{proof}
  This is the declared model definition of Circuit Sense.
\end{proof}
\begin{definition}[Slidewire Figure]
  \label{def:physics:phyx-mini-0919:phyxminiproblems-problemphyxmini0919-slidewirefigure}
  \lean{PhyXMiniProblems.ProblemPhyXMini0919.SlidewireFigure}
  \uses{def:physics:phyx-mini-0919:phyxminiproblems-problemphyxmini0919-conductorpart, def:physics:phyx-mini-0919:phyxminiproblems-problemphyxmini0919-spatialdirection, def:physics:phyx-mini-0919:phyxminiproblems-problemphyxmini0919-circuitsense}
  Literal content of image 919.png. The figure supplies directions and object labels but no numerical values for B, L, v, R, E, I, F, or P
\end{definition}
\begin{proof}
  This is the declared model definition of Slidewire Figure.
\end{proof}
\begin{definition}[Slidewire Circuit Setup]
  \label{def:physics:phyx-mini-0919:phyxminiproblems-problemphyxmini0919-slidewirecircuitsetup}
  \lean{PhyXMiniProblems.ProblemPhyXMini0919.SlidewireCircuitSetup}
  \uses{def:physics:phyx-mini-0919:phyxminiproblems-problemphyxmini0919-magneticfluxdensitymagnitude, def:physics:phyx-mini-0919:phyxminiproblems-problemphyxmini0919-lengthmagnitude, def:physics:phyx-mini-0919:phyxminiproblems-problemphyxmini0919-speedmagnitude, def:physics:phyx-mini-0919:phyxminiproblems-problemphyxmini0919-electromotiveforcemagnitude, def:physics:phyx-mini-0919:phyxminiproblems-problemphyxmini0919-electriccurrentmagnitude, def:physics:phyx-mini-0919:phyxminiproblems-problemphyxmini0919-resistancemagnitude, def:physics:phyx-mini-0919:phyxminiproblems-problemphyxmini0919-powermagnitude, def:physics:phyx-mini-0919:phyxminiproblems-problemphyxmini0919-forcemagnitude, def:physics:phyx-mini-0919:phyxminiproblems-problemphyxmini0919-slidewirefigure}
  Independent observables at one instant of the slidewire's motion. In particular, energyDissipationRate is not defined from the answer formula
\end{definition}
\begin{proof}
  This is the declared model definition of Slidewire Circuit Setup.
\end{proof}
\begin{definition}[Matches Slidewire Circuit Scenario]
  \label{def:physics:phyx-mini-0919:phyxminiproblems-problemphyxmini0919-matchesslidewirecircuitscenario}
  \lean{PhyXMiniProblems.ProblemPhyXMini0919.MatchesSlidewireCircuitScenario}
  \uses{def:physics:phyx-mini-0919:phyxminiproblems-problemphyxmini0919-slidewirecircuitsetup}
  The prose identifies the two displayed conductors as one closed circuit
\end{definition}
\begin{proof}
  This is the declared model definition of Matches Slidewire Circuit Scenario.
\end{proof}
\begin{definition}[Matches Primary Slidewire Figure]
  \label{def:physics:phyx-mini-0919:phyxminiproblems-problemphyxmini0919-matchesprimaryslidewirefigure}
  \lean{PhyXMiniProblems.ProblemPhyXMini0919.MatchesPrimarySlidewireFigure}
  \uses{def:physics:phyx-mini-0919:phyxminiproblems-problemphyxmini0919-slidewirecircuitsetup}
  Primary-raster evidence. These are qualitative direction readouts only; no power formula or answer-choice value occurs here
\end{definition}
\begin{proof}
  This is the declared model definition of Matches Primary Slidewire Figure.
\end{proof}
\begin{definition}[Has Physical Slidewire Parameters]
  \label{def:physics:phyx-mini-0919:phyxminiproblems-problemphyxmini0919-hasphysicalslidewireparameters}
  \lean{PhyXMiniProblems.ProblemPhyXMini0919.HasPhysicalSlidewireParameters}
  \uses{def:physics:phyx-mini-0919:phyxminiproblems-problemphyxmini0919-magneticfluxdensityinteslas, def:physics:phyx-mini-0919:phyxminiproblems-problemphyxmini0919-lengthinmeters, def:physics:phyx-mini-0919:phyxminiproblems-problemphyxmini0919-speedinmeterspersecond, def:physics:phyx-mini-0919:phyxminiproblems-problemphyxmini0919-resistanceinohms, def:physics:phyx-mini-0919:phyxminiproblems-problemphyxmini0919-slidewirecircuitsetup}
  Positivity and non-vacuity conditions for the physical slidewire state
\end{definition}
\begin{proof}
  This is the declared model definition of Has Physical Slidewire Parameters.
\end{proof}
\begin{definition}[Satisfies Circuit Resistance Composition]
  \label{def:physics:phyx-mini-0919:phyxminiproblems-problemphyxmini0919-satisfiescircuitresistancecomposition}
  \lean{PhyXMiniProblems.ProblemPhyXMini0919.SatisfiesCircuitResistanceComposition}
  \uses{def:physics:phyx-mini-0919:phyxminiproblems-problemphyxmini0919-resistanceinohms, def:physics:phyx-mini-0919:phyxminiproblems-problemphyxmini0919-slidewirecircuitsetup}
  The resistance named R in the problem is the sum of the resistance of the moving slidewire and that of the U-shaped conductor
\end{definition}
\begin{proof}
  This is the declared model definition of Satisfies Circuit Resistance Composition.
\end{proof}
\begin{definition}[Has Uniform Applied Magnetic Field]
  \label{def:physics:phyx-mini-0919:phyxminiproblems-problemphyxmini0919-hasuniformappliedmagneticfield}
  \lean{PhyXMiniProblems.ProblemPhyXMini0919.HasUniformAppliedMagneticField}
  \uses{def:physics:phyx-mini-0919:phyxminiproblems-problemphyxmini0919-magneticfluxdensityinteslas, def:physics:phyx-mini-0919:phyxminiproblems-problemphyxmini0919-slidewirecircuitsetup}
  The scalar flux-density magnitude calibrates the norm of the Physlib magnetic vector field throughout the region marked by uniform crosses
\end{definition}
\begin{proof}
  This is the declared model definition of Has Uniform Applied Magnetic Field.
\end{proof}
\begin{definition}[Satisfies Slidewire Induction And Circuit Laws]
  \label{def:physics:phyx-mini-0919:phyxminiproblems-problemphyxmini0919-satisfiesslidewireinductionandcircuitlaws}
  \lean{PhyXMiniProblems.ProblemPhyXMini0919.SatisfiesSlidewireInductionAndCircuitLaws}
  \uses{def:physics:phyx-mini-0919:phyxminiproblems-problemphyxmini0919-magneticfluxdensityinteslas, def:physics:phyx-mini-0919:phyxminiproblems-problemphyxmini0919-lengthinmeters, def:physics:phyx-mini-0919:phyxminiproblems-problemphyxmini0919-speedinmeterspersecond, def:physics:phyx-mini-0919:phyxminiproblems-problemphyxmini0919-electromotiveforceinvolts, def:physics:phyx-mini-0919:phyxminiproblems-problemphyxmini0919-electriccurrentinamperes, def:physics:phyx-mini-0919:phyxminiproblems-problemphyxmini0919-resistanceinohms, def:physics:phyx-mini-0919:phyxminiproblems-problemphyxmini0919-powerinwatts, def:physics:phyx-mini-0919:phyxminiproblems-problemphyxmini0919-forceinnewtons, def:physics:phyx-mini-0919:phyxminiproblems-problemphyxmini0919-slidewirecircuitsetup}
  School-physics laws for a rod perpendicular to both its velocity and a uniform magnetic field. These laws relate independent observables and do not state the requested expression for dissipated power
\end{definition}
\begin{proof}
  This is the declared model definition of Satisfies Slidewire Induction And Circuit Laws.
\end{proof}
% --- Archon declaration topology end ---
% --- Archon physics formalization source end ---
